% Generated draft: requires mathematical and source review.
\documentclass[11pt]{article}
\usepackage[margin=25mm]{geometry}
\usepackage{fontspec}
\setmainfont{Latin Modern Roman}
\usepackage{amsmath,amssymb,mathtools,unicode-math}
\setmathfont{Latin Modern Math}
\usepackage{enumitem,needspace,xurl,hyperref}
\setlength{\parindent}{0pt}
\setlength{\parskip}{5pt}
\setlength{\emergencystretch}{3em}
\newcommand{\N}{\mathbb{N}}
\newcommand{\Z}{\mathbb{Z}}
\newcommand{\Q}{\mathbb{Q}}
\newcommand{\R}{\mathbb{R}}
\newcommand{\C}{\mathbb{C}}
\newcommand{\F}{\mathbb{F}}
\newcommand{\A}{\mathbb{A}}
\newcommand{\PP}{\mathbb{P}}
\newcommand{\E}{\mathbb{E}}
\newcommand{\eps}{\varepsilon}
\newcommand{\dd}{\,\mathrm{d}}
\begin{document}
\section*{UnsolvedMath ID: 30003800}
\textbf{Source problem code:} catalogue entry
\subsection{Definitions and mathematical statement}
Let \(k\) be an algebraically closed field and let \(G\) be a semisimple algebraic group over \(k\). The source imposes no restriction on \(\operatorname{char}(k)\). The Dynkin type of \(G\) records the irreducible components of its root system. We interpret ``mixed Dynkin types'' provisionally as requiring at least two distinct irreducible Dynkin types among these components; the source does not explicitly define this term or restrict the central isogeny type of \(G\).

The notation \(BG\) denotes the classifying space of \(G\). The source writes \(F(BG)\), but does not define \(F\), its relationship to \(k\), or the construction of the classifying-space function field. One usual interpretation uses the invariant field \(k(V)^G\) of a generically free linear representation \(V\) of \(G\); this interpretation is not specified in the record.

For a field extension \(K/k\), once a coefficient object \(M\), a cohomology theory, and residue maps have been fixed, unramified degree-three cohomology consists of the classes in \(H^3(K,M)\) whose residues vanish at the specified valuations. For example, under the convention using all discrete rank-one valuations trivial on \(k\), write
\[
\mathcal V(K/k)=\{v: v\text{ is a discrete rank-one valuation of }K\text{ trivial on }k\}.
\]
When the relevant residue maps \(\partial_v\) are defined, this convention gives
\[
H^3_{\mathrm{nr}}(K/k;M)
=\bigcap_{v\in\mathcal V(K/k)}\ker\!\left(\partial_v:H^3(K,M)\longrightarrow\text{the residue group at }v\right).
\]
This explains the terminology without selecting coefficients or a residue theory for the source question; these choices, including their meaning in positive characteristic, remain unresolved.

The supplied question is whether, for every algebraically closed field \(k\) and every semisimple \(k\)-group \(G\) of mixed Dynkin types,
\[
H_{\mathrm{nr}}^3(F(BG))=0.
\]
Here triviality means that every unramified degree-three class is zero. This is a statement-only formulation of the supplied target. It is not a fully specified cohomological assertion until the meaning of \(F(BG)\), the coefficients, the residue convention, and the intended mixed types are fixed.
\subsection{Short English statement}
The question asks whether every unramified degree-three cohomology class of the classifying-space function field of a semisimple group with mixed Dynkin types vanishes over an algebraically closed field. The source leaves the function-field construction, coefficients, and characteristic-dependent conventions unspecified. Its open-status label is supplied curation, not an independently verified conclusion.
\subsection{Sources}
\begin{itemize}
\item Local UnsolvedMath registry: \texttt{data/unsolvedmath/problems.json}, permanent ID \(30003800\), source problem code \texttt{catalogue entry}, title ``Unramified Cohomology of Mixed-Type Classifying Spaces.'' The supplied record belongs to the Oberwolfach Reports Open Problems collection, gives proposed year \(2018\), and names no individual proposer.
\item Supplied report citation: \emph{Quadratic Forms and Related Structures over Fields} (2019), Oberwolfach Reports, pp.~1235--1291, DOI \texttt{10.4171/owr/2018/21}. \url{https://doi.org/10.4171/owr/2018/21}. The background describes Oberwolfach Reports as documenting mathematical workshops and their open-problem discussions. No exact page or contributor for this question is supplied.
\item Supplied literature reference, labeled ``Baek on unramified cohomology'': \url{https://doi.org/10.1016/j.jpaa.2021.106718}. The record supplies no full title or theorem identifying the applicable result.
\item Supplied literature reference, labeled ``Oberwolfach Report 21'': \url{https://doi.org/10.4171/owr/2018/21}. Literature checked at: \texttt{2026-08-21}. This is supplied dated curation, not independently verified. Its status is ``open'' and its assessment says that known results settle homogeneous exceptional and other restricted Dynkin types, but not the extracted question for arbitrary mixed semisimple Dynkin types. The background explicitly describes this as dated literature triage derived from the supplied OWR curation and calls for independent verification of time-sensitive status and cited claims.
\end{itemize}
\end{document}
